\documentclass{article}
\usepackage{natbib}
\usepackage{amsmath}
\usepackage{amsthm}
\usepackage{amssymb}
\usepackage{mathtools}
\usepackage{algorithm}
\usepackage{algpseudocode}
\usepackage{xcolor}

\usepackage[autostyle]{csquotes}

\usepackage{enumitem}
\newenvironment{QandA}{\begin{enumerate}[label=\bfseries\alph*.]\bfseries}
                      {\end{enumerate}}
\newenvironment{answered}{\par\normalfont}{}
\usepackage{lipsum}
\usepackage{hyperref}
\pagestyle{empty}

% ---- Change-highlight macros (were used but never defined in the original) ----
\newcommand{\rtwo}[1]{\textcolor{blue}{#1}}      % Reviewer 2 changes
\newcommand{\rthree}[1]{\textcolor{teal}{#1}}    % Reviewer 3 changes
\newcommand{\audit}[1]{\textcolor{magenta}{#1}}  % consistency/audit fixes

\usepackage[
    backend=biber,
    style=ieee,
  ]{biblatex}

\addbibresource{lagrangian.bib} %Imports bibliography file
\renewcommand*{\multicitedelim}{} %citations with no spaces in groups

\newtheorem{theorem}{Theorem}
\newtheorem{lemma}[theorem]{Lemma}
\newtheorem{proposition}[theorem]{Proposition}
\newtheorem{definition}{Definition}
\newtheorem{remark}{Remark}
\newtheorem{assumption}{Assumption}
\newtheorem{corollary}[theorem]{Corollary}

\title{Response to Reviewers}

\author{Anonymous Author(s)}

\begin{document}
\maketitle

Re: JCOMP-D-25-01237R1

\medskip

We thank the Associate Editor and both reviewers for the careful re-reading of our manuscript and for the constructive feedback. We particularly appreciate Reviewer~2's observation that our previous response letter contained claims that were not fully reflected in the manuscript itself.
Taking this criticism seriously, we have performed a \emph{full consistency pass} between this response letter and the revised manuscript before resubmission. Every change described below carries a corresponding section reference, and we have audited the manuscript end-to-end (including all tables, figures, and cross-references) to ensure no further discrepancies remain.

In the revised manuscript, we have indicated changes by color:
{\color{blue}\textbf{blue}} for changes responding to Reviewer~2,
{\color{teal}\textbf{teal}} for changes responding to Reviewer~3, and
{\color{magenta}\textbf{magenta}} for consistency and notation fixes that
address concerns raised across both reviewers (and that we identified during
the audit).

\section*{Summary of major revisions}

\begin{enumerate}[leftmargin=*,itemsep=0.3em]
    \item All occurrences of ``1D Navier--Stokes'' have been corrected to
    ``1D Burgers equation'' throughout the manuscript and table captions, since
    the governing equation $\partial_t u + u\partial_x u = \nu\partial_{xx} u$
    used in our 1D experiments is the viscous Burgers equation. (The only
    remaining textual occurrence was in a commented-out, non-compiled
    paragraph, which has now been deleted from the source.)
    \item Symbol overloading ($\lambda$, $w$, $c$) has been resolved through a
    systematic notation pass.
    \item Theorem~1 has been relabeled Proposition~1, made fully self-contained
    by an explicit definition paragraph, and its proof sketch moved to the
    appendix. \audit{The heuristic crossover result (former Theorem~3) has
    likewise been reclassified as a Proposition; see A.3.7.}
    \item A roadmap paragraph has been added at the end of \S1.
    \item Figures~2, 4, and~7 have been reformatted for readability, and the
    Figure~2 caption has been corrected.
    \item \audit{The experimental figures have been fully regenerated so that
    every convergence figure and weight-evolution figure is derived from the
    same run logs as the corresponding table (see A.1.1 and A.3.8).}
    \item A numerical audit of Tables~2--5 identified several inconsistencies,
    which have been corrected (see \S\ref{sec:audit} below).
\end{enumerate}

\begin{QandA}
Reviewer #2: The revised manuscript is improved, but I do not think it is yet ready for acceptance. The authors have clearly made an effort to expand the discussion and revise parts of the theory and experiments, but several important concerns remain only partially addressed. There are still noticeable inconsistencies between the response letter and the revised manuscript itself. The response states that some low-value theorem material was removed and that terminology was corrected throughout, yet related issues remain visible in the manuscript (for instance, the proof of Theorem 1 does not appear to add substantial value and could be moved to the appendix, and the revised manuscript still contains occurrences of 1D Navier-Stokes). I recommend that the authors perform a full consistency pass between the response letter and the actual manuscript before any further consideration.
\end{QandA}

\begin{answered}
We thank the reviewer for this pointed observation; it was justified, and we
have addressed both specific items, together with a full consistency pass:

\textbf{(a) Proof of Theorem~1 moved to the appendix.} The proof sketch
formerly appearing in the main text has been moved to
Appendix~A.2.1 in the revised manuscript, leaving only the proposition statement and the \emph{Link to Experimental Results} discussion (which we believe earns its place in the main narrative) in \S4.4. A one-line forward reference (``The derivation is given in Appendix~A.2.1.'') replaces the in-text proof.

\textbf{(b) ``1D Navier--Stokes'' corrected to ``1D Burgers''.} The reviewer
is correct: the governing equation we use in 1D,
$\partial_t u + u\partial_x u = \nu\partial_{xx} u$, is the viscous Burgers
equation, not Navier--Stokes. We have globally
replaced ``1D Navier--Stokes'' with ``1D Burgers'' (or ``1D viscous Burgers''
where the viscosity is salient). Mixed phrases such as ``1D and 2D
Navier--Stokes'' have been rewritten to ``1D Burgers and 2D Navier--Stokes.''
\audit{We further note that the single remaining occurrence of ``1D Navier
Stokes'' was located inside a commented-out (non-compiled) paragraph in the
source; that paragraph has now been deleted, so no occurrence remains in
either the compiled PDF or the source.}

\textbf{(c) Full consistency pass.} Beyond these two items, we have audited
the manuscript for further discrepancies between the response letter and the
text. The remaining points raised by Reviewer~2 below are addressed
individually. Additional inconsistencies we identified during the audit
are listed in \S\ref{sec:audit}.
\end{answered}

\begin{QandA}
A major remaining issue is notation. The paper still overloads symbols,
especially $\lambda$, which is used for source weights while also appearing
in standard augmented-Lagrangian notation as a multiplier in the background
discussion. The manuscript also shifts between $\lambda$ and $w$ for weights,
and uses $c$ both for a cost vector and for generic constraint functions.
These choices make the theory harder to follow than necessary and should be
cleaned up systematically.
\end{QandA}

\begin{answered}
We agree, and have made the following systematic changes:

\textbf{(i) Resolving the $\lambda$ collision.}
In the background sections (\S2.3, \S2.4) we previously wrote standard
augmented-Lagrangian and ADMM formulas with $\lambda$ as the multiplier,
which collides with our use of $\lambda$ as the simplex source-weight vector
throughout \S3--\S6. We have renamed the generic multiplier in the background
formulas from $\lambda$ to $y$. Equation~(3) now reads
\[
    \mathcal{L}_\rho(x, y) = f(x) + y^\top c(x) + \frac{\rho}{2}\|c(x)\|^2,
\]
and Eqs.~(4)--(7) have been updated accordingly. A note on first use
explains the choice:
``We use $y$ for the generic Lagrange multiplier in this background section to reserve $\lambda$ for source weights used throughout the rest of the paper.''

\textbf{(ii) Eliminating the $\lambda$/$w$ shift.}
Source weights are now denoted $\lambda$ \emph{everywhere}, including the
Data Assimilation background/analysis update (\S5.6), with $K$ reserved for
the Kalman gain.

\textbf{(iii) Reserving $c$ for the cost vector.}
The symbol $c$ is now used exclusively for the per-source cost vector
$c \in \mathbb{R}^n_+$. As part of this pass we also
corrected a mathematical error in the proof sketch of (former) Theorem~1:
because the budget constraint enters the augmented Lagrangian through the
one-sided penalty $\tfrac{\rho}{2}\max\{0,\,c^\top\lambda - B\}^2$,
its gradient with respect to $\lambda$ is $\rho\,\max\{0,\,c^\top\lambda - B\}\,c$,
not the ill-defined expression previously written. The corrected expression
appears in Proposition~1 (\S4.4) and in the moved proof (Appendix~A.2.1).
\end{answered}

\begin{QandA}
The presentation of Theorem 1 is also still not sufficiently clear. The theorem compresses several method-specific gradient forms into one statement, but several symbols are not defined clearly enough for the result to be self-contained.
\end{QandA}

\begin{answered}
We have rewritten the statement to be fully self-contained, and re-classified
it. Specifically:

\textbf{(i) Definition paragraph preceding the statement.}
A new paragraph immediately preceding the result defines every symbol
(the Softmax logits $\alpha$ and Jacobian $J_\sigma$; the multipliers
$\mu,\nu,\nu_{\text{budget}}$; the penalty $\rho$; the cost vector $c$ and
budget $B$; the tangent-space projection $\Pi_{\mathcal{T}_\Delta}$; and the
ADMM auxiliary variable $z_{k+1}$ with the simplex indicator subdifferential
$\partial I_{\Delta^{n-1}}$).

\textbf{(ii) Each row of the equation now has a correct, method-specific LHS}
(gradient w.r.t.\ logits for Softmax; a projected update direction for the
single-timescale variant; the augmented-Lagrangian gradient with the corrected
one-sided budget term; and a first-order optimality \emph{inclusion} for ADMM).

\textbf{(iii) Reclassification as Proposition~1.}
The result derives method-specific update rules from a unified Lagrangian via
routine calculus; we agree this is better labeled as a proposition. It is now
\rtwo{Proposition~1 (Method-Specific Weight Update Directions)}. \audit{Of the
remaining numbered results, Theorems~2 and~4 retain their designation, while
the crossover result (\S4.7) is now stated as Proposition~3, since it is a
heuristic order-of-magnitude scaling estimate rather than a rigorously proven
bound (see A.3.7).}

\textbf{(iv) Proof moved to appendix} (per Point~2.1).
\end{answered}

\begin{QandA}
The introduction would also benefit from a short roadmap paragraph at the end. At present, it moves from the contribution bullets directly into the next section, which is not ideal for a paper of this length and complexity.
\end{QandA}

\begin{answered}
We have added a roadmap paragraph immediately following the three contribution
bullets in \S1:

\emph{``The remainder of the paper is organized as follows.
Section~2 surveys related work in multi-task and multi-modal learning,
mixture-of-experts routing, constrained deep learning, multi-timescale
optimization, physical-systems integration, data assimilation, and the
machine-learning architectures we use. Section~3 introduces the augmented
Lagrangian formulation with budget constraints. Section~4 develops the
unified theoretical framework: notation and basic assumptions (\S4.1--\S4.2),
the unified problem (\S4.3), the method-specific weight update directions
(Proposition~1, \S4.4), convergence rates (Theorem~2, \S4.5--\S4.6), the
ADMM/Augmented-Lagrangian crossover analysis (\audit{Proposition~3}, \S4.7), and
numerical precision bounds (Theorem~4, \S4.8). Section~5 reports four
experimental studies---expert routing, multi-resolution learning, large-scale
source integration, and data assimilation---across 1D Burgers and 2D
Navier--Stokes regimes. Section~6 discusses trade-offs between Softmax and
constrained methods, and Section~7 concludes with future directions.''}
\end{answered}

\begin{QandA}
Finally, figure presentation still needs improvement. While the figures appear to have been revised, readability remains an issue, and the Figure 2 caption appears to contain formatting or rendering problems.
\end{QandA}

\begin{answered}
We have reformatted the figures and corrected the caption:

\textbf{(i) Readability.} All convergence and weight-evolution figures have
been regenerated at higher resolution with substantially enlarged fonts
(titles, axis labels, ticks, and legends), tightened inter-panel spacing, and
thicker lines, and are exported as vector PDF so they remain crisp at any zoom.
The stacked weight-evolution panels now fill the full plot width and use
descriptive legends (architecture names for expert routing; resolutions for
multi-resolution) rather than generic ``Source~$i$'' labels.

\textbf{(ii) Figure~2 caption.} The revised caption now identifies each expert:
\rtwo{``Experts: $\lambda_1$ (FiniteDiff), $\lambda_2$ (CNN), $\lambda_3$ (FNO),
$\lambda_4$ (DeepONet).''} and the en-dash rendering artifact
(``$\lambda_1$--$\lambda_4$'') has been fixed.
\end{answered}

\begin{QandA}
Reviewer 3: The manuscript establishes a theoretical framework for optimal weight learning in multi-source scientific machine learning [\ldots]. After carefully reading the revised manuscript, I find that the authors have clearly and satisfactorily addressed the previously raised concerns [\ldots]. I believe the paper is suitable for publication after careful reformatting and final polishing.
\end{QandA}

\begin{answered}
    We thank the reviewer for the positive assessment and have carried out the
    requested final reformatting and polishing (notation, figure readability,
    table formatting, and the consistency audit in \S\ref{sec:audit}).
\end{answered}

\section*{Additional consistency audit}
\label{sec:audit}

In addition to the points raised by the reviewers, we performed a full
numerical and cross-reference audit of the revised manuscript and identified
several remaining inconsistencies, all of which have now been corrected. We
list them here in the interest of transparency and to support the editor's
required ``changes made'' documentation.

\subsection*{A.1 Table corrections}

\begin{enumerate}[leftmargin=*,itemsep=0.4em]
    \item \audit{Table~3 (multi-resolution) fully re-extracted from the run
    logs.} Several 1D rows previously contained weights that did not sum to
    unity, together with Test-MSE and cost values that were not consistent
    with the corresponding convergence figures. The entire table (Test MSE,
    cost, wall-clock time, and $\lambda_1,\lambda_2,\lambda_3$ weights, plus the
    single-resolution baselines) has been re-extracted from the same run logs
    that generate Figures~3--4, so that table, figures, and log summaries now
    agree exactly. In the 1D block the lowest Test~MSE is achieved by ADMM
    ($3.90\times10^{-5}$), with the Improved and standard Lagrangian close
    behind ($5.30\times10^{-5}$ and $5.40\times10^{-5}$); all weight rows now
    sum to unity up to rounding, and the corresponding text in \S5.4 has been
    updated to match.

    \item \audit{Table~3, 1D, baseline costs.} The single-resolution baseline
    costs now follow the cost-doubling pattern
    (Base~32: $1.000$, Base~64: $2.000$, Base~128: $4.000$) consistently in
    both the 1D and 2D blocks.

    \item \audit{Source count in \S5.5.} The Problem Formulation, Eq.~(30), and
    surrounding text now consistently use $N = 100$ sources (matching Table~4's
    caption ``selection of 10 sources out of 100'' and the architecture
    specifications in Tables~B.10--B.11).

    \item \audit{Timing columns in Tables~2 and~3 (2D rows).}
    The 2D timing values in Tables~2 and~3 were previously near-identical due to
    a copy-paste error, despite reflecting different experiments. The Table~3
    timing column has been re-extracted from the multi-resolution run logs and
    now differs from Table~2 (e.g., 2D Softmax: $2052.97$\,s in Table~2 vs
    $3729.00$\,s in Table~3). All timing columns were re-checked against the
    logs.

    \item \audit{Time-column formatting.} All Time~(s) columns now use a uniform
    two-decimal-seconds format without thousands separators.

    \item \audit{Method ordering across tables.} All tables now follow a uniform
    order: Softmax, Lagrangian, ImpLagrangian, AugLag, ADMM.

    \item \audit{Method-name standardization.} ``Improved Augmented Lagrangian''
    is written in full on first occurrence in each section and abbreviated
    ``ImpLagrangian'' consistently thereafter; ``Auglag''/``AugLag'' are unified
    as ``AugLag'' throughout.

    \item \audit{Bold-best convention.} For every table claiming ``best results
    in bold,'' we re-verified that the bolded entry is the column optimum among
    the routing methods; several previously mis-bolded entries (e.g., in the
    multi-resolution 1D block) have been corrected.

    \item \audit{Consistent metric reported across convergence figures.} We
    verified that every convergence figure (Figs.~1, 3, 5, 6) plots the
    \emph{test} MSE/loss, matching the Test-MSE columns of Tables~2--5; the
    plotting scripts were corrected where they had previously read the
    train/validation column.
\end{enumerate}

\subsection*{A.2 In-text inconsistencies}

\begin{enumerate}[leftmargin=*,itemsep=0.4em]
    \item \audit{Conclusion, $\lambda_3$ value.} The Conclusion now refers to
    the AugLag value $\lambda_3 = 0.734$ (2D expert routing), consistent with
    Table~2; the spurious value $0.660$ has been removed.

    \item \audit{Duplicate paragraph in \S4.8.} The near-duplicate
    \emph{Link to Experimental Results} paragraph (which incorrectly cited
    ``$\lambda_i \approx 0.044$ in Table~5'') has been removed; the surviving
    paragraph correctly cites ``$S_1 \approx 0.035$ in Table~5'' (1D Softmax).

    \item \audit{Binary-convergence claim qualified.} The claim that Lagrangian
    weights converge to exact $0/1$ values has been qualified to hold where it
    is supported: exact $0/1$ in the 1D Data Assimilation regime (Table~5), and
    sparse boundary-attaining distributions in 2D (e.g.\ $\lambda=(0.537,0.463,0,0)$).

    \item \audit{Unresolved cross-reference.} The stray ``Fig.~??'' reference in
    \S5.5 has been fixed.

    \item \audit{DeepONet sensor count.} \S5.3 now specifies ``128 sensors in 1D
    and 4096 sensors (the full $64\times 64$ grid) in 2D,'' consistent with the
    fact that the expert-routing 2D experiment operates on a $64\times64$ grid
    (Table~B.9).

    \item \audit{Appendix-table captions.} The doubled ``learning'' in the
    captions of Tables~B.10--B.11 has been removed, and the ``Multi-Scale'' vs
    ``Multi-resolution'' caption wording has been unified.

    \item \audit{Data-assimilation budget stated.} The resource budget used in
    the Data Assimilation study is now stated explicitly in \S5.6
    ($B=2.0$ in 1D and $B=2.5$ in 2D), matching the final cost column of
    Table~5 (1D costs $\approx 2.0$; 2D costs $\le 2.5$) and the run
    configuration.
\end{enumerate}

\subsection*{A.3 Mathematical, experimental, and structural corrections}

\begin{enumerate}[leftmargin=*,itemsep=0.4em]
    \item \audit{Theorem~4, Softmax row.} The proof sketch now identifies the
    correct mechanism for the $\kappa\to\infty$ blow-up: it arises from
    $\kappa_\alpha$ (the condition number of the logit-to-weight map), which
    diverges as $\min_i\lambda_i\to 0$. The statement is unchanged.
    \audit{We additionally clarified that the bounded quantity is the
    error in the preconditioned Adam \emph{update} direction, with the raw
    gradient error obeying $\|\hat\nabla L_m-\nabla L_m\|\le\epsilon_{\text{machine}}\kappa_m\|\nabla L_m\|$
    (Assumption~3), amplified by the Adam preconditioner.}

    \item \audit{Theorem~2 redundancy.} The two previously identical bounds for
    the Augmented and Improved Augmented Lagrangian are now a single equation
    with method-dependent constants, noting
    $C_\lambda^{\text{ImpLag}} \le C_\lambda^{\text{AugLag}}$. \audit{We also
    corrected the ADMM primal rate to $O(1/\sqrt{k})$ (consistent with the
    non-convex coupling), with only the budget feasibility decaying
    geometrically.}

    \item \audit{Algorithm~1 projection symbol.} Algorithm~1 and Proposition~1
    now consistently use the tangent-space projection $\Pi_{\mathcal{T}_\Delta}$.

    \item \audit{Erroneous citation removed.} The unrelated RGB-T saliency
    reference formerly cited for multi-timescale ADMM has been replaced with
    \cite{zeng2021admm_constrained_problems_convergence_saturation, chen2017note_admm_constrained_problems, boyd2011distributed}.
    \audit{A duplicate bibliography entry for the same multi-fidelity review
    (previously cited under two keys, appearing as two separate references) has
    also been de-duplicated, and the ``adaptive ensemble methods'' citation in
    \S1 repointed to an ensemble-methods reference.}

    \item \audit{Abstract claim re Single-timescale AL.} The abstract now reads:
    \audit{``\ldots~For single-timescale Augmented Lagrangian coupling, by
    contrast, we establish sufficient conditions under which it achieves
    sublinear convergence governed by an effective condition number that
    quantifies the curvature imbalance between primal and dual variables.''}

    \item \audit{Section-overview ordering in \S5.} The \S5 overview paragraph
    now matches the actual subsection order.

    \item \audit{Reclassification of the crossover result and proofs added.}
    The Adaptive Crossover Point (former Theorem~3) is now stated as
    Proposition~3 and explicitly described as a heuristic scaling estimate,
    since its exact form (in particular the Adam factor $(1-\beta_1)/(1-\beta_2)$)
    is an order-of-magnitude argument rather than a proven bound; the
    per-iteration cost model is grounded in the sort-based $O(n\log n)$ simplex
    projection actually used in the code. In compensation we added full proofs
    of Theorems~2 and~4 in the appendix (descent-lemma/two-time-scale and
    floating-point error-propagation arguments, respectively).

    \item \audit{Grid sizes reconciled with the released code.} The large-scale
    integration study (\S5.5) now correctly states a 64-point grid in 1D and a
    $64\times64$ grid in 2D, matching the implementation; the previous
    ``256-point / $128\times128$'' description has been corrected. The data
    assimilation study remains at 128 points (1D) and $128\times128$ (2D) as in
    the code.

    \item \audit{Learning-rate schedule reconciled with the code (two-time-scale).}
    To make the experiments faithful to the two-time-scale analysis, we updated
    the code so that the source weights are optimized on a slower time scale than
    the model parameters: primal $\eta_\theta=10^{-5}$ for the model parameters and
    $\eta_\lambda=10^{-7}$ for the source weights (i.e., $\eta_\lambda \ll \eta_\theta$,
    realizing $\eta_\lambda=O(\eta_\theta^2)$ of Assumption~4), dual/multiplier
    ascent $10^{-2}$--$10^{-4}$, and weight decay $10^{-7}$, with both primal step
    sizes following a step-decay schedule (halved every $150$ epochs). \S5.2 now
    states these learning rates and \S4.5/\S5.2 describe the two-time-scale,
    diminishing-step regime under which the asymptotic $O(1/\sqrt{k})$/$O(1/k)$
    rates of Theorem~2 apply.
\end{enumerate}

\bigskip
\hrule
\bigskip

\section*{Data availability}

\begin{answered}
\rthree{All data used in our study were generated using the accompanying code,
which will be made available at
\url{https://github.com/ppokhrel1/multiple_lambdas_optimization}.}
\end{answered}

\bigskip

\noindent We are grateful to both reviewers for their careful attention to
this work, and to Reviewer~2 in particular for the rigor with which the
consistency between the response letter and the manuscript was checked. We
hope the present revision and accompanying audit make the manuscript ready
for acceptance.

\bigskip

\noindent With best regards,

\medskip

\noindent Pujan Pokhrel, on behalf of the authors

\printbibliography

\end{document}
